{\bfseries

Deutsch-Jozsa con la función paridad. Considere la función $f:\{0,1\}^3 \rightarrow \{0,1\}$ definida por $f(x) = x_1 \oplus x_2 \oplus x_3$, es decir, la paridad de los tres bits.
Aplique el algoritmo de Deutsch-Jozsa paso a paso para $n=3$, mostrando el estado después de cada etapa (inicialización, primera Hadamard, oráculo, segunda Hadamard).
Calcule explícitamente el estado final y verifique que la probabilidad de medir $\vert 000 \rangle$ en el primer registro es cero, lo que indica que la función es balanceada. (Puede usar la fórmula general o hacer el cálculo término a término.)}

\vspace{.5cm}
\begin{quote}

\begin{enumerate}
    \item \textbf{Inicialización }
        $$
        \ket{\psi_0} = \ket{000} \otimes \ket{1}
        $$

    \item \textbf{Primera Hadamard sobre el primer registro.*} 
        \begin{align*}
            \ket{\psi_1} = H^{\otimes 4} \ket{\psi_{0}} &= \frac{1}{\sqrt{2^4}} (\ket{0}+\ket{1})
            \otimes (\ket{0}+\ket{1})\otimes (\ket{0}+\ket{1}) \otimes (\ket{0}-\ket{1}) \\
                                                        &= \frac{1}{\sqrt{2^3}} \left(\sum _{x
                                                                \in \{0,1\}^3} \ket{x} \right)
                                                                \otimes \ket{-}
        \end{align*}

    \item \textbf{Aplicación del oráculo $U_f$}

        Otra vez vamos a aplicar la regla del phase kickback:
        \begin{align*}
            \ket{\psi_2} = U_f \ket{\psi_1} &= \frac{1}{\sqrt{2^3}} U_f \left(\sum _{x
                                                                \in \{0,1\}^3} \ket{x} \right)
                                                                \otimes \ket{-} \\ 
                                            &= \frac{1}{\sqrt{2^3}} (\ket{000} - \ket{001} - \ket{010} +
                                            \ket{011} - \ket{100} + \ket{101} + \ket{110} -
                                            \ket{111} )\ket{-}
        \end{align*}

    \item \textbf{Segunda Hadamard sobre el primer registro.}

        Antes que nada voy a factorizar para que sea mas facil, quitamos el factor de normalizacion
        y el ultimo registro por ahora:
        \begin{align*}
            \ket{0} (\ket{00}-\ket{01}-\ket{10}+\ket{11}) - \ket{1} (\ket{00}
            -\ket{01}-\ket{10}+\ket{11}) &= (\ket{0}-\ket{1}) \otimes
            (\ket{00}-\ket{01}-\ket{10}+\ket{11}) \\ 
            (\ket{0}-\ket{1}) \otimes ( \ket{0} (\ket{0}-\ket{1}) - \ket{1}(\ket{0}-\ket{1})) &=
            (\ket{0}-\ket{1}) \otimes (\ket{0}-\ket{1})\otimes (\ket{0}-\ket{1}) 
        \end{align*}

        Ahora si:
        \begin{align*}
            \ket{\psi_3} = H^{\otimes 3} \ket{\psi_2} = H^{\otimes 3} \ket{-} \otimes \ket{-}
            \otimes \ket{-} \otimes \ket{-} = \ket{111} \otimes \ket{-} 
        \end{align*}
\end{enumerate}

    \textbf{Calcular la probabilidad de medir $\ket{000}$ en el primer registro y confirmar que la
    función es balanceada.}

    $$
    P(000) = |\bra{000}\ket{111}| ^2 = 0
    $$

    Al ser ortogonales la probabilidad es 0 y demostramos que es una función balanceada. 
\end{quote}
\vspace{.5cm}


